% Section 10 — Conclusion (compressed)
\section{Conclusion}
\label{sec:conclusion}

Kafka's exactly-once semantics is a precise guarantee about Kafka's own durable state, and it
stops at the external-store boundary. For a double-entry ledger that boundary is not an
implementation detail: a duplicate posting is a silent accounting error, and a lost posting is a
silent omission that no retry can distinguish from a new event. This paper made the boundary
explicit, enumerated the crash points at which it manifests, and turned the question ``is this
pipeline exactly once?'' into a set of falsifiable invariant-preservation claims per mechanism
and per failure class, backed by a machine-checked bounded model.

The evaluation in Section~\ref{sec:results} is populated from frozen run records by the
pre-registered analysis; its hypothesis verdicts and decision table map each failure exposure to
the minimal mechanism set that covers it and the measured cost of each guarantee. Two conclusions
follow for practice. First, ``we use Kafka transactions'' is not a statement about the ledger:
the mechanism must be described by which durable transitions share a site and which idempotency
or progress-atomicity argument closes the remaining window. Second, no mechanism is free --- the
guarantees that survive a crash between effect and progress cost throughput or latency, and the
decision table makes that cost explicit rather than hiding it behind a single ``exactly once''
label.

Two directions remain open: extending the ledger model to multi-leg and multi-currency entries,
reversals and corrections (which multiply the crash points), and moving from a single host to a
replicated multi-node deployment with real machine failures, where the durability assumptions the
model abstracts become observable. The artifact, injectors and oracle are public so both
extensions can build on the measured baseline.
